\clearpage
\section{Results and Discussion}

The final system satisfies the \textbf{core goals} of the assignment while remaining
small and easy to reason about.

\subsection{Observed outcome}

The project successfully distributes the alarm system over multiple Pekko
Cluster nodes without changing its domain rules. Sensors, keypad, and control
logic can run on different JVMs, yet they still interact only through typed
actor messages.

The most relevant design outcome is that the \textbf{alarm state remains centralized}
even in the clustered version. This avoids the inconsistency that would arise
if multiple independent control units were active at the same time.

\subsection{Discussion of trade-offs}

The solution intentionally stops at the level required by the assignment.
There is no persistence layer, no replicated state, and no attempt to keep the
alarm state after a crash. This is a deliberate limitation rather than a
missing feature, because the specification explicitly allows state loss after
recreation as long as the system returns to a safe recovery mode.

From a \textbf{didactic point of view}, this trade-off is reasonable. The code shows:

\begin{itemize}
    \item cluster formation through seed nodes;
    \item role-based node startup;
    \item remote service discovery;
    \item message serialization;
    \item failure-aware recovery.
\end{itemize}

At the same time, it avoids turning the assignment into a persistence or
replication project, which would shift the focus away from the intended
learning goals.
